\documentclass[a4paper,12pt]{article}
\usepackage[utf8]{inputenc}
\usepackage[T2A]{fontenc}
\usepackage[english, russian]{babel}

\usepackage{csquotes}
\usepackage{epigraph}
\usepackage{tcolorbox}

% Определяем окружение для Формата ввода
\newtcolorbox{inputformat}[1][]{colback=blue!5!white, colframe=blue!65!black,
    fonttitle=\bfseries, title=Формат ввода, #1}

% Определяем окружение для Формата вывода
\newtcolorbox{outputformat}[1][]{colback=green!5!white, colframe=green!65!black,
    fonttitle=\bfseries, title=Формат вывода, #1}

% Определяем окружение для комментария
\newtcolorbox{exercisecomment}[1][]{colback=yellow!5!white, colframe=yellow!80!black, fonttitle=\bfseries, title=Комментарий, #1}

% Определяем окружение для замечания
\newtcolorbox{exercisenote}[1][]{colback=red!5!white, colframe=red!80!black, fonttitle=\bfseries, title=Замечание, #1}

\usepackage{geometry}
\geometry{left=2cm}
\geometry{right=2cm}
\geometry{top=2cm}
\geometry{bottom=2cm}

\usepackage{amsmath, amsfonts, amsthm, amssymb, amsopn, amscd}
\usepackage{enumerate}
\usepackage{enumitem}
\usepackage[mathscr]{eucal}

\usepackage{hyperref}
\hypersetup{unicode=true,final=true,colorlinks=true}

\theoremstyle{remark}
\newtheorem{exercise}{\textbf{Упражнение}}[section]
\renewcommand{\theexercise}{\textbf{\arabic{exercise}}}

\title{Листок 07. Условный оператор в Python}
\author{А.Н. Лата}

\begin{document}

\section*{\centering Листок 07. Условный оператор в Python}

\begin{exercisenote}[title=Замечания]
\begin{itemize}
    \item Для решения приведённых ниже упражнений не требуется создавать (определять) пользовательские функции.
    \item Данные вводятся с клавиатуры с помощью функции {\color{blue}{input()}}. Не забывайте приводить входные данные к нужному типу.
    \item При решении упражнений полученные результаты выводите на экран с помощью функции {\color{blue}{print()}}.
    \item Позаботьтесь о том, чтобы выводимый на экран результат был снабжён информацией о нём, там где это необходимо.
    \item Не забывайте писать комментарии к вашему коду.
\end{itemize}
\end{exercisenote}

\textbf{Упражнения}
\begin{enumerate}

    % 1
    \item Дано число. Проверьте, положительное оно, отрицательное или ноль.

    % 2
    \item Дано целое число. Проверьте, является ли число чётным.
    \textit{Если число чётное, вывести \enquote{Чётное}, иначе — \enquote{Нечётное}.}

    % 3
    \item Даны два числа. Вывести большее из них.
    \textit{Если два числа равны, вывести \enquote{Числа равны}, иначе — вывести большее из них.}

    % 4
    \item Дано целое число. Определите, является ли число трёхзначным.

    % 5
    \item Дано целое число. Проверьте, кратно ли оно 7.

    % 6
    \item Дано целое число. Определите, делится ли оно на 2 и не делится на 3 одновременно.

    % 7 — уточнено: не вычислять произведение
    \item Даны три целых числа $a$, $b$, $c$. Определите знак их произведения, \textbf{не вычисляя само произведение}: проанализируйте знаки каждого из чисел с помощью условных операторов.

    % 8
    \item Дано число. Проверьте, лежит ли число в интервале от $-5$ до $3$ (не включая границы).
    \textit{Если число больше $-5$ и меньше $3$, вывести \enquote{Число в диапазоне}, иначе — \enquote{Число вне диапазона}.}

    % 9
    \item Дано целое число. Определите, является ли число отрицательным и чётным одновременно.
    \textit{Если число отрицательное и чётное, вывести \enquote{Отрицательное чётное}; если отрицательное и нечётное — \enquote{Отрицательное нечётное}; если неотрицательное — \enquote{Неотрицательное}.}

    % % 10
    % \item Дана строка. Проверьте, начинается ли строка с русской буквы \enquote{А} (заглавной).

    % % 11
    % \item Дана строка. Проверьте, является ли первая буква строки заглавной.

    % % 12
    % \item Дана строка. Определите, оканчивается ли строка на точку.

    % % 13
    % \item Дана строка. Проверьте, является ли строка пустой.

    % % 14
    % \item Дана строка. Проверьте, содержит ли строка пробелы.

    % % 15
    % \item Дана строка. Проверьте, состоит ли строка только из цифр.
    % \textit{Подсказка: воспользуйтесь методом строк \texttt{.isdigit()}.}

    % 16 — уточнено условие
    \item Дана строка, состоящая только из строчных русских букв без пробелов.
    Проверьте, является ли строка палиндромом.
    \textit{Палиндром — строка, которая читается одинаково слева направо и справа налево (например, \enquote{шалаш}, \enquote{казак}). Вывести \enquote{Палиндром}, если условие выполнено, иначе — \enquote{Не палиндром}.}

    % 17
    \item Дан список. Определите, равны ли первый и последний элементы списка.

    % 18 — заменено на более содержательное
    \item Дан список целых чисел. Проверьте, все ли элементы списка положительны.
    \textit{Вывести \enquote{Все положительные}, если условие выполнено, иначе — \enquote{Есть неположительные}.
    Подсказка: воспользуйтесь функцией \texttt{min()}.}

    % 19
    \item Дан кортеж, содержащий числа. Проверьте, содержит ли кортеж значение 7.
    \textit{Вывести \enquote{Семёрка найдена}, иначе — \enquote{Семёрка отсутствует}.}

    % 20
    \item Дан список из пяти чисел. Определите, является ли элемент с индексом~2 чётным числом.
    \textit{Если элемент на индексе 2 чётный, вывести \enquote{Чётное}, иначе — \enquote{Нечётное}.}

    % 21
    \item Определите, является ли прямоугольник квадратом. Стороны прямоугольника вводятся пользователем.

    % 22
    \item Определите, лежит ли точка с координатами $(x, y)$ в первой четверти координатной плоскости. \textit{Точка лежит в первой четверти, если $x > 0$ и $y > 0$.}

    % 23 — добавлена подсказка
    \item Проверьте, может ли существовать треугольник с тремя заданными сторонами $a$, $b$, $c$.
    \textit{Треугольник существует, если сумма любых двух его сторон больше третьей, то есть должны выполняться три условия: $a + b > c$, $a + c > b$, $b + c > a$.}

    % 24 — уточнён граничный случай
    \item Проверьте, лежит ли точка \textbf{строго внутри} круга радиуса $R$ с центром в начале координат. \textit{Точка $(x, y)$ находится строго внутри круга, если $x^2 + y^2 < R^2$. Точка на окружности не считается лежащей внутри круга.}

    % 25
    \item Проверьте, является ли введённый пользователем символ гласной буквой русского алфавита.
    \textit{Гласные русского алфавита: а, е, ё, и, й, о, у, ы, э, ю, я (и их заглавные варианты).}

    % 26
    % \item Проверьте, состоит ли строка только из заглавных букв.
    % \textit{Подсказка: воспользуйтесь методом строк \texttt{.isupper()}.}

    % 27
    \item Определите, лежит ли точка $(x, y)$ строго внутри прямоугольника. Прямоугольник задан координатами двух противоположных углов $(x_1, y_1)$ и $(x_2, y_2)$, где $x_1 < x_2$ и $y_1 < y_2$.

    % 28 — добавлена подсказка и уточнение типа данных
    \item Проверьте, является ли треугольник прямоугольным. Стороны треугольника $a$, $b$, $c$ — \textbf{целые числа}.
    \textit{Треугольник является прямоугольным, если сумма квадратов двух меньших сторон равна квадрату наибольшей. Подсказка: воспользуйтесь функцией \texttt{sorted()}, чтобы упорядочить стороны.}

    % 29 — добавлена подсказка
    \item Проверьте, образуют ли три числа арифметическую прогрессию.
    \textit{Три числа образуют арифметическую прогрессию, если разности между соседними элементами в упорядоченном ряду равны. Подсказка: воспользуйтесь функцией \texttt{sorted()}, чтобы упорядочить числа.}

    % 30 — добавлены примеры
    \item Проверьте, является ли число \textbf{автоморфным} — то есть его квадрат оканчивается этим же числом.
    \textit{Например: $5^2 = 25$ — оканчивается на 5, значит 5 — автоморфное; $76^2 = 5776$ — оканчивается на 76, значит 76 — автоморфное; $3^2 = 9$ — не оканчивается на 3, значит 3 — не автоморфное.
    Подсказка: переведите число и его квадрат в строку с помощью \texttt{str()}.}

    % 31 — добавлена подсказка
    \item Проверьте, все ли цифры введённого пользователем натурального числа различны.
    \textit{Подсказка: переведите число в строку с помощью \texttt{str()} и воспользуйтесь функцией \texttt{set()}.}

    % 32 — добавлена подсказка
    \item Проверьте, является ли сумма первой и последней цифры введённого пользователем натурального числа чётной.
    \textit{Подсказка: переведите число в строку с помощью \texttt{str()} и обращайтесь к цифрам как к символам строки.}

    % 33
    \item Дано положительное четырёхзначное число. Проверьте, равна ли сумма первой и последней цифры числа произведению средних цифр.
    \textit{Подсказка: переведите число в строку с помощью \texttt{str()}.}

    % 34 — новая задача со словарём
    \item Дан словарь с оценками студента, где ключи — названия предметов, значения — оценки (целые числа от 2 до 5). Проверьте, есть ли в словаре предмет \enquote{Математика}. Если есть — выведите оценку; если нет — выведите \enquote{Предмет не найден}. Дополнительно: если оценка равна 2, вывести \enquote{Задолженность}.

\end{enumerate}

\section*{\centering Усложнённые задачи}
\textit{Для каждой задачи используйте вложенные условные операторы \texttt{if-else} для обработки всех возможных комбинаций входных данных. Важно учесть все граничные случаи и корректно обработать некорректные входные данные.}

\begin{enumerate}
    \setcounter{enumi}{34}

    % 35
    \item Напишите программу, которая принимает возраст человека и определяет, к какой возрастной категории он относится:
    \begin{itemize}
        \item 0--2: младенец
        \item 3--6: дошкольник
        \item 7--17: школьник
        \item 18--22: студент
        \item 23--60: работник
        \item 61 и старше: пенсионер
    \end{itemize}
    При этом возраст не может быть отрицательным: если введено отрицательное число, выводите сообщение \enquote{Некорректный возраст}.

\end{enumerate}

\setcounter{exercise}{35}
\begin{exercise}[Стоимость билета в кинотеатр]
Создайте программу определения стоимости билета в кинотеатр. Учитывайте:
\begin{itemize}
    \item Время сеанса (до 12:00, 12:00--16:00, после 16:00)
    \item День недели (будни/выходные)
    \item Возраст посетителя (дети до 7 лет, школьники 7--17, взрослые 18+)
    \item Наличие студенческого билета
\end{itemize}
Базовая цена: 500 рублей. \\
\textbf{Формула расчёта:}
$$\text{Итоговая цена} = \text{Базовая цена} \times K_{\text{время}} \times K_{\text{день}} \times K_{\text{возраст}}$$

\noindent Где: \\[4pt]
\textit{Временной коэффициент $K_{\text{время}}$:}
\begin{itemize}
    \item До 12:00: $\times 0{,}8$
    \item 12:00--16:00: $\times 0{,}9$
    \item После 16:00: $\times 1{,}0$
\end{itemize}

\textit{Коэффициент дня $K_{\text{день}}$:}
\begin{itemize}
    \item Выходной: $\times 1{,}2$
    \item Будний: $\times 1{,}0$
\end{itemize}

\textit{Возрастной коэффициент $K_{\text{возраст}}$:}
\begin{itemize}
    \item До 7 лет: $\times 0{,}5$
    \item 7--17 лет: $\times 0{,}7$
    \item 18+ со студенческим: $\times 0{,}8$
    \item 18+ без студенческого: $\times 1{,}0$
\end{itemize}

\begin{inputformat}[title=Входные данные]
    \texttt{time} (int): время сеанса (часы, 0--23) \\
    \texttt{is\_weekend} (bool): выходной день \\
    \texttt{age} (int): возраст посетителя \\
    \texttt{has\_student\_card} (bool): наличие студенческого билета
\end{inputformat}

\begin{outputformat}[title=Пример 1]
    Входные данные: \\
    \texttt{time = 10} (до 12:00) \\
    \texttt{is\_weekend = False} \\
    \texttt{age = 15} \\
    \texttt{has\_student\_card = False} \\[4pt]
    Расчёт: \\
    1. Базовая цена: 500 руб. \\
    2. Временной коэффициент: $\times 0{,}8$ (утро) \\
    3. День недели: $\times 1{,}0$ (будни) \\
    4. Возрастной коэффициент: $\times 0{,}7$ (школьник) \\[4pt]
    $500 \times 0{,}8 \times 1{,}0 \times 0{,}7 = 280$ руб.
\end{outputformat}

\begin{outputformat}[title=Пример 2]
    Входные данные: \\
    \texttt{time = 19} (после 16:00) \\
    \texttt{is\_weekend = True} \\
    \texttt{age = 20} \\
    \texttt{has\_student\_card = True} \\[4pt]
    Расчёт: \\
    1. Базовая цена: 500 руб. \\
    2. Временной коэффициент: $\times 1{,}0$ (вечер) \\
    3. День недели: $\times 1{,}2$ (выходной) \\
    4. Возрастной коэффициент: $\times 0{,}8$ (студент) \\[4pt]
    $500 \times 1{,}0 \times 1{,}2 \times 0{,}8 = 480$ руб.
\end{outputformat}

\end{exercise}

\begin{exercise}[Расчёт стоимости доставки]
Создайте программу расчёта стоимости доставки посылки. \\
Параметры:
\begin{itemize}
    \item Вес (до 1 кг, 1--5 кг, 5--10 кг, более 10 кг)
    \item Расстояние (до 5 км, 5--20 км, более 20 км)
    \item Срочность (обычная/экспресс)
    \item Хрупкость груза (да/нет)
    \item Тип доставки (дверь-дверь / до пункта выдачи)
\end{itemize}

\noindent\textbf{Формула расчёта:}

\begin{multline*}
\text{Итоговая цена} = \bigl(\text{Базовая цена} + \text{Надбавка за вес} + \text{Надбавка за расстояние}\bigr) \times K_{\text{срочн.}} \times K_{\text{хрупк.}} + \\ + \text{Надбавка за доставку до двери}
\end{multline*}

\noindent Где: \\[4pt]
Базовая цена $= 300$ руб. \\[4pt]

\textit{Надбавка за вес:}
\begin{itemize}
    \item До 1 кг: $+100$ руб.
    \item 1--5 кг: $+200$ руб.
    \item 5--10 кг: $+300$ руб.
    \item Более 10 кг: $+500$ руб.
\end{itemize}

\textit{Надбавка за расстояние:}
\begin{itemize}
    \item До 5 км: $+100$ руб.
    \item 5--20 км: $+200$ руб.
    \item Более 20 км: $+300$ руб.
\end{itemize}

\textit{Коэффициент срочности $K_{\text{срочн.}}$:}
\begin{itemize}
    \item Экспресс: $\times 1{,}5$
    \item Обычная: $\times 1{,}0$
\end{itemize}

\textit{Коэффициент хрупкости $K_{\text{хрупк.}}$:}
\begin{itemize}
    \item Хрупкий: $\times 1{,}3$
    \item Обычный: $\times 1{,}0$
\end{itemize}

Надбавка за доставку до двери: $+200$ руб.

\begin{inputformat}[title=Входные данные]
    \texttt{weight} (float): вес в кг \\
    \texttt{distance} (float): расстояние в км \\
    \texttt{is\_express} (bool): экспресс-доставка \\
    \texttt{is\_fragile} (bool): хрупкий груз \\
    \texttt{door\_delivery} (bool): доставка до двери
\end{inputformat}

\begin{outputformat}[title=Пример 1]
    Входные данные: \\
    \texttt{weight = 0.5} кг \\
    \texttt{distance = 3} км \\
    \texttt{is\_express = False} \\
    \texttt{is\_fragile = False} \\
    \texttt{door\_delivery = True} \\[4pt]
    Расчёт: \\
    1. Базовая цена: 300 руб. \\
    2. Надбавка за вес: $+100$ руб. (до 1 кг) \\
    3. Надбавка за расстояние: $+100$ руб. (до 5 км) \\
    4. Коэффициент срочности: $\times 1{,}0$ (обычная) \\
    5. Коэффициент хрупкости: $\times 1{,}0$ (обычный) \\
    6. Надбавка за доставку до двери: $+200$ руб. \\[4pt]
    $(300 + 100 + 100) \times 1{,}0 \times 1{,}0 + 200 = 700$ руб.
\end{outputformat}

\begin{outputformat}[title=Пример 2]
    Входные данные: \\
    \texttt{weight = 3} кг \\
    \texttt{distance = 12} км \\
    \texttt{is\_express = True} \\
    \texttt{is\_fragile = True} \\
    \texttt{door\_delivery = False} \\[4pt]
    Расчёт: \\
    1. Базовая цена: 300 руб. \\
    2. Надбавка за вес: $+200$ руб. (1--5 кг) \\
    3. Надбавка за расстояние: $+200$ руб. (5--20 км) \\
    4. Коэффициент срочности: $\times 1{,}5$ (экспресс) \\
    5. Коэффициент хрупкости: $\times 1{,}3$ (хрупкий) \\
    6. Надбавка за доставку до двери: $+0$ руб. \\[4pt]
    $(300 + 200 + 200) \times 1{,}5 \times 1{,}3 + 0 = 1365$ руб.
\end{outputformat}

\end{exercise}

\end{document}